\chapter{Mar.~17 --- Sheaf Cohomology, Part 3}

\section{Examples of Sheaf Cohomology}

\begin{remark}
  We want to
  compute $H^i(\PP^n, \OO_{\PP^n}(m))$
  and use it to give applications.
\end{remark}

\begin{example}
  We have already seen
  in Example \ref{ex:pp1} that
  \[
    H^i(\PP^1, \OO_{\PP^1})
    =
    \begin{cases}
      k & \text{if } i = 0, \\
      0 & \text{if } i > 0,
    \end{cases}
  \]
  and that (note that we have
  $\Omega_{\PP^1} \cong \OO_{\PP^1}(-2)$)
  \[
    H^i(\PP^1, \OO_{\PP^1}(-2))
    = \begin{cases}
      k & \text{if } i = 1, \\
      0 & \text{otherwise}.
    \end{cases}
  \]
\end{example}

\begin{example}
  We have seen in homework that
  $H^0(\PP^n, \OO_{\PP^n}(m)) = \Gamma(\PP^n, \OO_{\PP^n}(m)) = k[x_0, \dots, x_n]_m$.
\end{example}

\begin{example}
  Let $X = \PP^n$,
  $\mathcal{U} = \{U_i\}_{i = 0}^n$ for
  $U_i = D(x_i)$, and
  $S = \bigoplus_{m \ge 0} S_m = \bigoplus_{m \ge 0} k[x_0, \dots, x_n]_m$.
  We want to compute $H^i$ of the complex
  \[
    \begin{tikzcd}
      \prod_i \OO_{\PP^n}(m)(U_i)
      \ar[r] & \prod_{i < j} \OO_{\PP^n}(m)(U_{i, j})
      \ar[r] & \cdots \ar[r] &
      \OO_{\PP^n}(m)(U_{0, \dots, n})
    \end{tikzcd}
  \]
  We know that for $U_{i_0, \dots, i_r} = D(x_{i_0} \cdots x_{i_r})$,
  we have
  \[
    \OO_{\PP^n}(m)(U_{i_0, \dots, i_r})
    = (S_{x_{i_0} \cdots x_{i_r}})_m.
  \]
  To make this simpler, write
  $\mathcal{F} = \bigoplus_{m \in \Z} \OO_X(m)$.
  Now
  \[
    H^i(X, \mathcal{F})
    = \bigoplus_{m \in \Z} H^i(X, \OO_X(m)),
  \]
  and
  observe that $\mathcal{F}(U_{i_0, \dots, i_r}) = S_{x_{i_0} \cdots x_{i_r}}$.
  So we have
  \[
    C^\bullet(\mathcal{U}, \mathcal{F})
    = \left[
      \prod_i S_{x_i}
      \longrightarrow
      \prod_{i < j} S_{x_i x_j}
      \longrightarrow \cdots
      \longrightarrow
      S_{x_0 \cdots x_n}
    \right].
  \]
  Then we want to compute
  $H^i(C^\bullet(\mathcal{U}, \mathcal{F}))$
  and keep track of the grading. For
  $i = 0$,
  \[
    H^0(X, \mathcal{F})
    = \ker\bigg(
      \prod_i S_{x_i} \longrightarrow \prod_{i < j} S_{x_i x_j}\bigg)
      = S,
  \]
  so $H^0(\PP^n, \OO_{\PP^n}(m)) = S_m$.
  In particular,
  $H^0(\PP^n, \OO_{\PP^n}(m)) = 0$ for $m < 0$.
  For $i = n$,
  \[
    H^n(X, \mathcal{F})
    = \coker\bigg(
      \prod_i S_{x_0 \cdots \widehat{x}_i \cdots x_n}
      \longrightarrow
      S_{x_0 \cdots x_n}
    \bigg)
    = \bigoplus_{a_i < 0} k \overline{x_0^{a_0} \cdots x_n^{a_n}}.
  \]
  So $H^n(X, \OO_{\PP^n}(m)) = \bigoplus_{a_i < 0, \sum a_i = m} k \overline{x_0^{a_0} \cdots x_n^{a_n}}$.
  In particular,
  $H^n(X, \OO_{\PP^n}(m)) = 0$ for
  $m \ge -n$ and
  \[
    H^n(X, \OO_{\PP^n}(-n - 1))
    = k \overline{x_0^{-1} \cdots x_n^{-1}}.
  \]
  Observe that there is a perfect pairing
  \begin{align*}
    H^0(X, \OO_{\PP^n}(m))
    \times H^n(X, \OO_{\PP^n}(-m - n - 1))
    &\longrightarrow H^n(X, \OO_{\PP^n}(-n - 1)) \cong k \\
    (x_0^{a_0} \cdots x_n^{a_n}, \overline{x_0^{b_0} \cdots x_n^{b_n}})
    &\longmapsto
    \overline{x_0^{a_0 + b_0} \cdots x_n^{a_n + b_n}}.
  \end{align*}
  where $a_i \ge 0$ and
  $\sum a_i = m$, and
  $b_i < 0$ and $\sum b_i = -m - n - 1$.
  Note that
  \[
    \overline{x_0^{a_0 + b_0} \cdots x_n^{a_n + b_n}}
    =
    \begin{cases}
      \overline{x_0^{-1} \cdots x_n^{-1}} & \text{if } a_i + b_i = -1 \text{ for all } i, \\
      0 & \text{otherwise}.
    \end{cases}
  \]
  This says that
  $H^n(X, \OO_{\PP^n}(-m) \otimes \omega_{\PP^n}) \cong H^0(X, \OO_{\PP^n}(m))^\vee$,
  which is a special case of \emph{Serre duality}.

  Now for the middle cohomology, we
  claim that $H^i(X, \OO_{\PP^n}(m)) = 0$
  for $0 < i < n$.
  For this, it suffices to show that
  $H^i(X, \mathcal{F}) = 0$ for $0 < i < n$.
  We argue by induction.
  Consider the short exact sequence
  \[
    \begin{tikzcd}
      0 \ar[r]
      & S(-1) \ar[r, "\times x_n"]
      & S \ar[r]
      & S / x_n S \ar[r] & 0
    \end{tikzcd}
  \]
  Then taking $M \mapsto \widetilde{M}$
  gives a short exact sequence
  \[
    \begin{tikzcd}
    0 \ar[r] & \OO_X(-1)
    \ar[r] & \OO_X \ar[r] &
    i_* \OO_H \ar[r] & 0
    \end{tikzcd}
  \]
  with $H = V(x_n) \cong \PP^{n - 1} \hookrightarrow \PP^n$.
  Set $\mathcal{F}_H = i_* (\bigoplus_{m \in \Z} \OO_H(m))$.
  Applying $- \otimes \mathcal{F}$
  gives
  \[
    \begin{tikzcd}
      0 \ar[r] & \mathcal{F}(-1)
      \ar[r, "\times x_n"]
     & \mathcal{F} \ar[r]
      & \mathcal{F}_H \ar[r] & 0
    \end{tikzcd}
  \]
  This gives a long exact sequence
  \[
    \begin{tikzcd}
    H^0(X, \mathcal{F}(-1))
    \ar[r, "\times x_n", hook] & H^0(X, \mathcal{F})
    \ar[r, two heads] & H^0(X, \mathcal{F}_H)
    \ar[r] & \cdots \\
    \ar[r] & H^i(X, \mathcal{F}(-1))
    \ar[r, "\times x_n"] & H^i(X, \mathcal{F})
    \ar[r] & H^i(X, \mathcal{F}_H)
    \ar[r] & \cdots \\
    \ar[r] & H^{n - 1}(X, \mathcal{F}(H))
    \ar[r, "\delta", hook] & H^n(X, \mathcal{F}(-1))
    \ar[r, "\times x_n"] & H^n(X, \mathcal{F})
    \ar[r] & H^n(X, \mathcal{F}_H)
    \end{tikzcd}
  \]
  Note that $H^i(X, \mathcal{F}_H) = 0$
  for $1 \le i < n - 1$ by induction.
  One can check that $\delta$ is
  multiplication by $1 / x_n$, so it
  is injective.
  Thus for $0 < i < n$,
  we have an isomorphism
  (note that $\mathcal{F}(-1) \cong \mathcal{F}$)
  \[
    H^i(X, \mathcal{F}(-1))
    \overset{\times x_n}{\longrightarrow}
    H^i(X, \mathcal{F}).
  \]
  So $H^i(X, \mathcal{F})$, viewed
  as an $S$-module, satisfies that
  $x_n : H^i(X, \mathcal{F}) \to H^i(X, \mathcal{F})$ is an isomorphism.
  Now
  \[
    H^i(X, \mathcal{F})_{x_n}
    \cong H^i(C^\bullet(\mathcal{U}, \mathcal{F}))_{x_n}
    \cong H^i(C^\bullet(\{U_i \cap U_n\}_{i = 0}^n, \mathcal{F}|_{U_n})) = 0
  \]
  as $\{U_i \cap U_n\}_{i = 0}^n$ is
  an open cover of $U_n$, and $U_n$ is
  an affine variety. So
  $H^i(X, \mathcal{F}) = 0$ for $0 < i < n$.
  This then implies that
  $H^i(\PP^n, \OO_{\PP^n}(m)) = 0$ for
  $0 < i < n$. From the duality perspective,
  we have
  \[
    H^i(X, \OO_{\PP^n}(m))
    \cong H^{n - i}(X, \OO_{\PP^n}(m)^\vee \otimes \omega_{\PP^n})^\vee.
  \]
\end{example}

\section{Applications to Coherent Sheaves}

\begin{remark}
  For this section, let
  $X \hookrightarrow \PP^n$ be a projective
  variety and $\mathcal{F}$ a
  quasicoherent sheaf on $X$.
\end{remark}

\begin{prop}
  $H^i(X, \mathcal{F}) = 0$ for
  $i > \dim X$.
\end{prop}

\begin{proof}
  Let $r = \dim X$, and choose
  hyperplanes $H_0, \dots, H_r$ on
  $X$ such that
  $X \cap (H_0 \cap \cdots \cap H_r) = \varnothing$.
  So
  \[
    X = U_0 \cup \cdots U_r
  \]
  with $U_i = X \setminus H_i$. As the
  $U_i$ are affine, we get
  $H^i(X, \mathcal{F})
    = H^i(C^\bullet(\{U_i\}_{i = 0}^r, \mathcal{F}))
    = 0$
  for $i > r$.
\end{proof}

\begin{prop}
  If $\mathcal{F}$ is coherent,
  then $H^i(X, \mathcal{F})$ is a
  finite dimensional $k$-vector space.
\end{prop}

\begin{proof}
  See the reduction given in Section
  \ref{sec:motivation}, where the only
  missing part was that
  $H^i(\PP^n, \OO_{\PP^n}(m))$
  is finite dimensional for $i \ge 0$ and
  $m \in \Z$.
\end{proof}

\begin{theorem}[Serre vanishing]\label{thm:serre-vanishing}
  If $\mathcal{F}$ is a coherent,
  there exists $m_0 = m_0(\mathcal{F})$
  such that
  \[
    H^i(X, \mathcal{F}(m))
    = 0
  \]
  for $i > 0$ and $m \ge m_0$.
\end{theorem}

\begin{proof}
  Note that $i_*(\mathcal{F}(m)) \cong (i_* \mathcal{F})(m)$, and
  we have
  \[
    H^i(X, \mathcal{F}(m))
    \cong H^i(\PP^n, i_* \mathcal{F}(m)),
  \]
  so we may reduce to the case
  $X = \PP^n$. Then there is a
  short exact sequence of coherent sheaves
  \[
    \begin{tikzcd}
      0 \ar[r] & \mathcal{E} \ar[r] &
      \bigoplus_{i = 1}^r \OO_X(a_i)
      \ar[r] & \mathcal{F} \ar[r] & 0
    \end{tikzcd}
  \]
  Twisting by $\OO_X(m)$ gives
  \[
    \begin{tikzcd}
      0 \ar[r] & \mathcal{E}(m) \ar[r] &
      \bigoplus_{i = 1}^r \OO_X(m + a_i)
      \ar[r] & \mathcal{F}(m) \ar[r] & 0
    \end{tikzcd}
  \]
  So we get a long exact sequence
  on cohomology
  \[
    \begin{tikzcd}
      H^i(\PP^n, \bigoplus_{i = 1}^r \OO_X(m + a_i))
      \ar[r] & H^i(\PP^n, \mathcal{F}(m)) \ar[r] &
      H^{i + 1}(\PP^n, \mathcal{E}(m))
      \ar[r] & \cdots
    \end{tikzcd}
  \]
  Now $H^{i + 1}(\PP^n, \mathcal{E}(m)) = 0$
  by induction and
  $H^i(\PP^n, \bigoplus_{i = 1}^r \OO_X(m + a_i)) = 0$
  for $m > -n - 1 - a_i$. So we see that
  $H^i(\PP^n, \mathcal{F}(m)) = 0$
  for $i > 0$ and $m \gg 0$.
\end{proof}

\begin{remark}
  The coherent assumption in
  Theorem \ref{thm:serre-vanishing} is
  necessary: Let $X = \PP^n$
  and $\mathcal{F} = \bigoplus_{m \in \Z} \OO_{\PP^n}(m)$.
\end{remark}

\begin{remark}
  We can get a cohomological
  characterization of ampleness. One
  can check that for an invertible sheaf
  $\mathcal{L}$ on a projective
  variety $X$, the following
  are equivalent:
  \begin{enumerate}
    \item $\mathcal{L}$ is ample;
    \item $\mathcal{F} \otimes \mathcal{L}^{\otimes m}$
      is globally generated for any
      coherent sheaf $\mathcal{F}$ and $m \gg 0$ (dependent on $\mathcal{F}$);
    \item $H^i(X, \mathcal{F} \otimes \mathcal{L}^{\otimes m}) = 0$
      for any coherent $\mathcal{F}$ on
      $X$, $i > 0$, and $m \gg 0$
      (dependent on $\mathcal{F}$).
  \end{enumerate}
  Note that in general (not necessarily
  for a projective variety), the
  definition of ampleness is given by (2).
\end{remark}
